% Theory slides (Assignment 8). Transcribed from Docs/Writing/manuscript/spec/spec.tex and spec/sections/sp1,sp2,sp3,sp5,sp6,sp7.
% GENERATED by .rc_tmp/a8/gen.py: statements copied from the spec; \ref and \eqref replaced by the spec's numbers;
% equation labels replaced by \tag{N} (spec equation numbers); enumerate options changed to beamer form. No proofs, no commentary.
\documentclass[aspectratio=169]{beamer}
\usetheme{Madrid}
\usepackage[T1]{fontenc}
\usepackage{lmodern}
\usepackage{amsmath,amssymb,mathtools,bm}
\usepackage{booktabs,array}
\usepackage{url}
\urlstyle{tt}
\DeclareUrlCommand\code{}
\setbeamertemplate{navigation symbols}{}
\setbeamertemplate{enumerate items}[default]
\setbeamersize{text margin left=7mm,text margin right=7mm}
\setlength{\emergencystretch}{3em}
\renewcommand{\arraystretch}{1.0}

% ---- math macros of Docs/Writing/shared/smcbvs.sty and of spec.tex (copied; \hyperref links dropped)
\newcommand{\defasm}[2]{\expandafter\def\csname asmdisp@#1\endcsname{#2}}
\newcommand{\asmdisp}[1]{\csname asmdisp@#1\endcsname}
\defasm{GL}{GL}\defasm{DIM}{DIM}\defasm{GR}{GR}\defasm{SP}{SP}\defasm{BC}{BC}\defasm{BM}{BM}
\defasm{FSb}{FS$_\beta$}\defasm{FSt}{FS$_\theta$}\defasm{EF}{EF}\defasm{NC}{NC}\defasm{EFS}{EF$_\Sigma$}
\defasm{V}{V}\defasm{NL}{NL}
\defasm{C1}{C1}\defasm{C2}{C2}\defasm{C2p}{C2$^+$}\defasm{C3}{C3}
\defasm{SGR}{SG}
\defasm{dist}{DIST}\defasm{reported}{RM}\defasm{margin}{MARG}\defasm{omitted}{OMIT}
\defasm{reported-fixed}{RMF}\defasm{SCR}{SCR}\defasm{WIT}{WIT}
\newcommand{\as}[1]{\textup{(\asmdisp{#1})}}
\newcommand{\R}{\mathbb{R}}
\newcommand{\N}{\mathbb{N}}
\newcommand{\one}{\mathbf{1}}
\newcommand{\ind}[1]{\mathbf{1}\{#1\}}
\newcommand{\set}[1]{[#1]}
\newcommand{\abs}[1]{\lvert #1\rvert}
\newcommand{\norm}[1]{\lVert #1\rVert}
\DeclareMathOperator{\E}{\mathbb{E}}
\renewcommand{\P}{\mathbb{P}}
\DeclareMathOperator{\Var}{Var}
\DeclareMathOperator{\Cov}{Cov}
\DeclareMathOperator{\Corr}{Corr}
\DeclareMathOperator{\sd}{sd}
\DeclareMathOperator{\supp}{supp}
\DeclareMathOperator*{\argmax}{arg\,max}
\newcommand{\Normal}{\mathrm{N}}
\newcommand{\IG}{\mathrm{IG}}
\newcommand{\BetaDist}{\mathrm{Beta}}
\newcommand{\BB}{\mathrm{BB}}
\newcommand{\OP}{O_{\P}}
\newcommand{\oP}{o_{\P}}
\newcommand{\lmin}{\lambda_{\min}}
\newcommand{\yr}{y^{\mathrm r}}
\newcommand{\Xr}{X^{\mathrm r}}
\newcommand{\praw}{p_{\mathrm{raw}}}
\newcommand{\sigs}{\sigma_*}
\newcommand{\bstar}{\beta^*}
\newcommand{\tstar}{\theta^*}
\newcommand{\tmin}{\theta^*_{\min}}
\newcommand{\Sz}{S_0}
\newcommand{\sz}{s_0}
\newcommand{\Rhat}{\hat R}
\newcommand{\Sig}{\Sigma_n}
\newcommand{\rn}{\sqrt{\log n/n}}
\newcommand{\SSq}{\mathrm{SS}}
\newcommand{\Hl}{\mathcal H}
\newcommand{\Dn}{\mathcal D_n}
\newcommand{\tlow}{\underline\theta}
\newcommand{\Om}{\Omega}
\newcommand{\pis}[1]{\pi(#1\mid y)}
\newcommand{\PIP}{\mathrm{PIP}}
\newcommand{\HPM}{\mathrm{HPM}}
\newcommand{\MPM}{\mathrm{MPM}}
\newcommand{\tSz}{\tilde S_0}
\newcommand{\Shat}{\hat S}
\newcommand{\Zhat}{\hat Z}
\newcommand{\PIPhat}{\widehat{\mathrm{PIP}}}
\newcommand{\MPMhat}{\widehat{\mathrm{MPM}}}
\newcommand{\Lat}{\mathcal L}
\newcommand{\Dom}{\mathfrak D}
\newcommand{\Acal}{\mathcal E_{\mathrm{cal}}}
\newcommand{\Nc}{N^{\mathrm c}}
\newcommand{\grid}[1]{\xi_{[#1]}}
\newcommand{\Lbi}{L_{\mathrm{bi}}}
\newcommand{\Ts}{T_{\mathrm s}}
\newcommand{\Tmax}{T^{\max}}
\newcommand{\rid}[1]{\texttt{#1}}

\title{SMC-BVS: assumptions and theorems}
\subtitle{proved and planned}
\author{SMC-BVS project, Writing Department}
\institute{}
\date{24 November 2026}

\begin{document}

\begin{frame}
\titlepage
\end{frame}

\section{Notation and setting}
\begin{frame}
\centering\Large\textbf{Notation and setting}
\end{frame}

\begin{frame}{Notation and setting (1): general notation}
\footnotesize
\begin{block}{Definition 1.1 (General notation)}
For an integer $k\ge1$, $\set{k}:=\{1,\dots,k\}$. Vectors are columns; $\norm{v}$ and $\norm{v}_1$ are the Euclidean and
$\ell_1$ norms; $I_n$ is the identity and $\one_n$ the vector of ones; $\lmin(\cdot)$ is the smallest eigenvalue of a
symmetric matrix; $\abs{S}$ is the size of a set; $a\vee b:=\max(a,b)$; $\ind{E}$ is an indicator. Logarithms are natural
($\log$) or binary ($\log_2$). A star marks a true value, and the superscript $\mathrm r$ marks raw, unstandardised data.
For a matrix $\Lambda$ and sets $U,V$, $\Lambda_{UV}$ is the sub-matrix with rows in $U$ and columns in $V$, and
$\Lambda_U:=\Lambda_{UU}$. For random variables, $V_n=\OP(a_n)$ means that $\abs{V_n}/a_n$ is bounded in probability and
$V_n=\oP(a_n)$ that it tends to zero in probability; \emph{with probability tending to one} means that the probability of
the event tends to one. For positive sequences, $O$, $o$, $\lesssim$, $\gtrsim$, $\asymp$ and $\gg$ have their usual
meaning ($a_n\gg b_n$: $b_n=o(a_n)$).
\end{block}
\end{frame}

\begin{frame}{Notation and setting (2): data and standardisation}
\footnotesize
\begin{block}{Definition 1.2 (Data and retained predictors)}
The data are a response $\yr\in\R^n$ and $\praw$ predictors. A predictor that is constant in the sample is
\emph{removed}; $\Xr$ is the $n\times p$ matrix of the $p$ retained predictors, and $p$ always denotes this number.
The retained predictors are indexed $1,\dots,p$ in the order of their original columns, and $\mathrm{col}(j)\in\set{\praw}$ is the
original one-based column number of retained predictor $j$ (the \emph{label map}). Models and all indices in this
document are retained indices; the method reports models, the median model and inclusion-probability labels in original
column numbers, $j\mapsto\mathrm{col}(j)$.
\end{block}
\begin{block}{Definition 1.4 (Standardisation)}
With $P^\perp_{\one}:=I_n-n^{-1}\one\one^\top$ and $D$ the diagonal matrix of the column sample standard deviations of
$\Xr$ (divisor $n-1$),
\begin{equation}\tag{1}
  X:=P^\perp_{\one}\Xr D^{-1},\qquad y:=P^\perp_{\one}\yr/s_y,\qquad s_y^2:=\norm{P^\perp_{\one}\yr}^2/(n-1).
\end{equation}
Every column $x_j$ of $X$ has mean zero and $\norm{x_j}^2=n-1$, and $y^\top y=n-1$.
\end{block}
\end{frame}

\begin{frame}{Notation and setting (3): models and truth}
\footnotesize
\begin{block}{Definition 1.5 (Models)}
A \emph{model} is a set $S\subseteq\set{p}$, equivalently its inclusion vector $\gamma\in\{0,1\}^p$ with
$\gamma_j:=\ind{j\in S}$; then $\abs{\gamma}=\abs{S}$ and $X_\gamma=X_S$. This document writes $S$. For a model $S$:
$X_S$ holds the columns of $X$ in $S$; $G_S:=X_S^\top X_S$ is their Gram matrix; $\Rhat_S:=G_S/(n-1)$ is their sample
correlation matrix; $\hat r_{jl}$ is the sample correlation of predictors $j$ and $l$; $P_S$ is the projection onto the
span of $X_S$ ($P_\emptyset:=0$); and $z_S:=X_S^\top y$.
\end{block}
\begin{block}{Definition 1.6 (True model and true coefficients)}
Under the data-generating model~\textup{(2)} (Definition~2.1), the \emph{true model} is
$\Sz:=\{j:\beta^*_j\neq0\}$, of size $\sz$; its members are the \emph{true predictors}, the others the \emph{null
predictors}. The \emph{standardised coefficients} are $\tstar:=D\bstar$ (support $\Sz$), so that
$P^\perp_{\one}\yr=X\tstar+P^\perp_{\one}\varepsilon$; the \emph{mean signal} is $\mu:=X\tstar$; and
$\tmin:=\min_{j\in\Sz}\abs{\theta^*_j}$.
\end{block}
\end{frame}

\begin{frame}{Notation and setting (4): designs and the distribution setting}
\footnotesize
\begin{block}{Definition 1.3 (Designs)}
In a \emph{fixed design}, $\Xr$ is deterministic and probability refers to the noise. In a \emph{random Gaussian design},
the rows of $\Xr$ are i.i.d.\ $\Normal_p(m_n,\Sig)$ and probability refers to the joint law of design and noise;
$W=(W_1,\dots,W_p)\sim\Normal_p(0,\Sig)$ denotes a centred copy of one row, with sub-vectors $W_S$. For sets $U$ and
indices $j,l$, $\Cov(W_j,W_l\mid W_U)$ is the partial covariance given $(W_k)_{k\in U}$, and
$\bar\beta_{U,l}:=\Sigma_{n,UU}^{-1}\Sigma_{n,Ul}$ is the population regression coefficient of $W_l$ on $W_U$, padded
with zeros to a vector indexed by a larger set when needed.
\end{block}
\begin{block}{Definition 1.8 (Setting of the distribution results)}
In Theorem~2 and the results that cite it, the observations are $n$ i.i.d.\ copies $(x^{(i)},y^{(i)})$ of a
pair $(x,y)$, with $x\in\R^{\praw}$ the vector of all observed predictors, and $\zeta_j:=(x_j-\E x_j)/\sd(x_j)$. The data
set is $\Dn$ and its law is $\P^*$. The method runs on the predictors that are not constant in the sample.
\end{block}
\end{frame}

\begin{frame}{Notation and setting (5): the model}
\footnotesize
\begin{block}{Definition 2.1 (Data-generating model)}
The data follow a Gaussian linear model with an intercept,
\begin{equation}\tag{2}
  \yr=\alpha^*\one_n+\Xr\bstar+\varepsilon,\qquad \varepsilon\sim\Normal_n(0,\sigs^2I_n),
\end{equation}
with $\sigs^2>0$ and $\varepsilon$ independent of $\Xr$. The theorems are proved under this model (in the forms fixed by
their assumptions), not under the working model of Definition~2.4.
\end{block}
\begin{block}{Definition 2.2 (Model space)}
The method considers every model with at most $K$ predictors,
\begin{equation}\tag{3}
  \Om:=\{S\subseteq\set{p}:\abs{S}\le K\},\qquad K:=\min\big(p,\lfloor n/2\rfloor\big).
\end{equation}
\end{block}
\end{frame}

\begin{frame}{Notation and setting (6): the working model}
\footnotesize
\begin{block}{Definition 2.3 (Parameters)}
The parameters of the working model are a model $S\in\Om$, coefficients $\theta_S\in\R^{\abs{S}}$ and a variance
$\sigma^2>0$. The intercept is removed by centring and has no parameter.
\end{block}
\begin{block}{Definition 2.4 (Working model and likelihood)}
Given $(S,\theta_S,\sigma^2)$, the method fits to the standardised data~\textup{(1)}
\[
  y\mid S,\theta_S,\sigma^2\ \sim\ \Normal_n(X_S\theta_S,\sigma^2I_n),
\]
with likelihood
\begin{equation}\tag{4}
  p(y\mid S,\theta_S,\sigma^2)=(2\pi\sigma^2)^{-n/2}\exp\Big\{-\frac{\norm{y-X_S\theta_S}^2}{2\sigma^2}\Big\}.
\end{equation}
\end{block}
\end{frame}

\begin{frame}{Notation and setting (7): the prior}
\footnotesize
\begin{block}{Definition 3.1 (Model prior, with the size cap $K$)}
The prior mass of a model $S\in\Om$ depends only on its size:
\begin{equation}\tag{5}
  \pi(S)=\frac{w(\abs{S})}{\bar w_K},\qquad w(k):=\frac{1}{2\binom{2p-1}{k}},\qquad
  \bar w_K:=\sum_{k=0}^{K}\binom{p}{k}w(k),
\end{equation}
and $\pi(S)=0$ for $\abs{S}>K$. This is the beta-binomial prior $\BB(1,p)$ restricted to $\Om$: $w(\abs{S})$ is the mass
of $S$ when $q\sim\BetaDist(1,p)$ and, given $q$, each predictor enters independently with probability $q$.
\end{block}
\end{frame}

\begin{frame}{Notation and setting (8): the prior (cont.)}
\footnotesize
\begin{block}{Definition 3.2 (Coefficient prior)}
Given $S$ and $\sigma^2$, the coefficients have the $g$-prior with $g=n$ and a unit ridge,
\begin{equation}\tag{6}
  \theta_S\mid S,\sigma^2\sim\Normal\big(0,\;n\sigma^2(G_S+I)^{-1}\big),
\end{equation}
with density
\[
  p(\theta_S\mid S,\sigma^2)=(2\pi n\sigma^2)^{-\abs{S}/2}\det(G_S+I)^{1/2}
  \exp\Big\{-\frac{\theta_S^\top(G_S+I)\theta_S}{2n\sigma^2}\Big\}.
\]
For $S=\emptyset$ there is no coefficient.
\end{block}
\begin{block}{Definition 3.3 (Variance prior)}
$\sigma^2\sim\IG(a_0,b_0)$ with $a_0=b_0=0.01$, independent of $S$, with density
\[
  p(\sigma^2)=\frac{b_0^{a_0}}{\Gamma(a_0)}(\sigma^2)^{-a_0-1}e^{-b_0/\sigma^2},\qquad \sigma^2>0 .
\]
\end{block}
\end{frame}

\begin{frame}{Notation and setting (9): the posterior}
\footnotesize
\begin{block}{Definition 3.4 (Joint posterior)}
The posterior of $(S,\theta_S,\sigma^2)$ is the law on $\bigcup_{S\in\Om}\{S\}\times\R^{\abs{S}}\times(0,\infty)$ with
density, with respect to counting measure on $\Om$ times Lebesgue measure,
\[
  \Pi(S,\theta_S,\sigma^2\mid y)\ \propto\ \pi(S)\,p(\theta_S\mid S,\sigma^2)\,p(\sigma^2)\,p(y\mid S,\theta_S,\sigma^2),
\]
with the factors of Definitions~3.1, 3.2, 3.3 and~2.4.
\end{block}
\begin{block}{Definition 3.6 (Log score, posterior over models, evidence)}
The \emph{log score} of $S\in\Om$, the \emph{posterior over models} and the \emph{evidence} are
\begin{equation}\tag{8}
  F(S):=\log\{w(\abs{S})/\bar w_K\}+\log m(S),\qquad \pis{S}=\frac{e^{F(S)}}{Z},\qquad Z:=\sum_{S'\in\Om}e^{F(S')} .
\end{equation}
\end{block}
\end{frame}

\begin{frame}{Notation and setting (10): posterior summaries}
\footnotesize
\begin{block}{Definition 3.7 (Posterior summaries)}
The \emph{inclusion probability} of predictor $j$ is $\PIP_j:=\sum_{S\ni j}\pis{S}$; the \emph{highest posterior model} is
$\HPM:=\argmax_S\pis{S}$; the \emph{median probability model} is $\MPM:=\{j:\PIP_j>1/2\}$; and the \emph{posterior miss}
is $1-\pis{\Sz}$.
\end{block}
\begin{block}{Definition 3.8 (Fitted conditional law and Hellinger distance)}
In the setting of Definition~1.8, a posterior draw $(S,\beta,\sigma^2)$, mapped back to the original
units, defines the conditional law $\Normal(\bar y+(x-\bar x)^\top\beta,\sigma^2)$ of $y$ given $x$; the law of $x$ is not
estimated. Let $f$ be the joint density of $(x,y)$ with $x$ from its true law and $y\mid x$ from this fitted law, and $f^*$
the true joint density. Then $d_H(f,f^*):=(1-\int\sqrt{ff^*})^{1/2}$, and $\Pi(\cdot\mid\Dn)$ denotes the posterior of
$(S,\beta,\sigma^2)$. The posterior is not defined if fewer than two predictors remain or the response is constant in the
sample.
\end{block}
\end{frame}

\begin{frame}{Notation and setting (11): the asymptotic regime}
\footnotesize
\begin{block}{Definition 1.7 (Asymptotic regime and rate parameters)}
Most results are asymptotic in $n\to\infty$, along a sequence of problems in which $p=p_n$ (and $\praw$), the design and
the coefficients may change with $n$. Lemma~S1, Theorem~6(a) and Proposition~5 hold for every data
set; Theorem~2 holds for every $n\ge n_0$; the delivery clauses hold for every admissible $n$. A quantity is
\emph{fixed} if it does not depend on $n$; every constant introduced
in an assumption is fixed. The parameters that may depend on $n$, and where they are introduced, are:
\end{block}
\end{frame}

\begin{frame}{Notation and setting (11): the asymptotic regime (cont.)}
\footnotesize
\begin{block}{Definition 1.7 (Asymptotic regime and rate parameters) (cont.)}
\begin{itemize}
  \item the dimensions $n$, $p=p_n$, $\praw$, and the true model size $\sz$ (fixed under \as{SP});
  \item the posterior rate $\rn$, the size cap $K=K_n$ (Definition~2.2), the block size cap $B=B_n$, the
    base move cap $\ell_n$ and the particle numbers $N_0$, $N_i$, $\Nc_i$, the burn-in $\Lbi$ and the step ceiling
    $\Tmax_i$ (Table~1);
  \item the level $M$ and the window $q_n$ of Assumption~\as{dist};
  \item the model-size budget $k_n$ of Assumption~\as{SCR};
  \item the sequences $\varpi_n$, $\delta_n$, $d_n$, $L_n$ and the leakage set $Q$ with neighbourhoods $\Hl_l$ and
    $N_l$ of Assumption~\as{reported};
  \item the sequence $r^\circ_n$ and neighbourhoods $\Hl_l$ of Assumption~\as{NL}; the sequence $\omega_n$ of
    Assumption~\as{reported-fixed}; and the rate $\rho_n$ of Remark~S15.
\end{itemize}
\end{block}
\end{frame}

\section{The assumptions}
\begin{frame}
\centering\Large\textbf{The assumptions}
\end{frame}

\begin{frame}{Assumptions: the data model}
\small
\begin{block}{Assumption 5.1 (\asmdisp{GL}) \textit{(Gaussian linear model)}}
The data follow~\textup{(2)}, with a fixed noise variance $\sigs^2>0$ and noise independent of the design.
\end{block}
\begin{block}{Assumption 5.2 (\asmdisp{DIM}) \textit{(Dimension)}}
$2\le p\le n^A$ for a fixed $A\ge1$.
\end{block}
\begin{block}{Assumption 5.3 (\asmdisp{SP}) \textit{(Sparsity)}}
The true model size $\sz\ge1$ is fixed.
\end{block}
\end{frame}

\begin{frame}{Assumptions: the predictors and the signal (1)}
\small
\begin{block}{Assumption 5.4 (\asmdisp{GR}) \textit{(Gaussian rows)}}
The rows of $\Xr$ are i.i.d.\ $\Normal_p(m_n,\Sig)$, where $\Sig$ is a correlation matrix.
\end{block}
\begin{block}{Assumption 5.5 (\asmdisp{SGR}) \textit{(Sub-Gaussian rows on sparse directions)}}
The raw rows are i.i.d.\ copies of a random vector $x\in\R^{\praw}$ with mean $m_n$ and covariance matrix $\Sig$, a
correlation matrix, and, for a fixed $L>0$, $\E\exp\{\lambda u^\top(x-m_n)\}\le\exp\{L^2\lambda^2u^\top\Sig u/2\}$ for every
$\lambda\in\R$ and every $u\in\R^{\praw}$ with at most $5(A+1)(\sz+1)\log n$ nonzero entries. Gaussian rows \as{GR} are
the case $L=1$.
\end{block}
\end{frame}

\begin{frame}{Assumptions: the predictors and the signal (2)}
\small
\begin{block}{Assumption 5.6 (\asmdisp{FSb}) \textit{(Fixed signals, raw scale)}}
The intercept $\alpha^*$ and the raw coefficients $\bstar$ are deterministic, and
$b_{\mathrm{lo}}\le\abs{\beta^*_j}\le b_{\mathrm{hi}}$ for every $j\in\Sz$, with fixed $0<b_{\mathrm{lo}}\le b_{\mathrm{hi}}<\infty$.
\end{block}
\begin{block}{Assumption 5.7 (\asmdisp{FSt}) \textit{(Fixed signals, standardised scale)}}
$\yr=\alpha\one+X\tstar+\varepsilon$ with $\varepsilon\sim\Normal_n(0,\sigs^2I_n)$, a fixed $\sigs^2>0$, noise independent
of $\Xr$, a deterministic $\tstar$, $\sz\ge1$ fixed, and $\theta_{\mathrm{lo}}\le\abs{\theta^*_j}\le\theta_{\mathrm{hi}}$ for
$j\in\Sz$, with fixed $0<\theta_{\mathrm{lo}}\le\theta_{\mathrm{hi}}$.
\end{block}
\begin{block}{Assumption 5.8 (\asmdisp{EFS}) \textit{(Population eigenvalue floor)}}
$\lmin(\Sigma_{n,S})\ge\varphi^2$ for every $S$ with $\abs{S}\le5(A+1)(\sz+1)\log n$, for a fixed $\varphi\in(0,1]$ (for
instance, $\lmin(\Sig)\ge\varphi^2$).
\end{block}
\end{frame}

\begin{frame}{Assumptions: the distribution-level setting}
\small
\noindent In the setting of Definition~1.8, fix a \emph{level} $M\ge1$; the \emph{window} is
$q_n:=\min\{\praw,\lceil C_qM\log n\rceil\}$ with $C_q:=2(s(A+2)+4)/\log2$.

\begin{block}{Assumption 5.9 (\asmdisp{dist}) \textit{(Distribution-level setting; A6 Assumption~3)}}
The integer $s\ge0$ and the constants $A\ge1$, $L\ge1$, $c_\psi>0$ and $\bar b>0$ are fixed.
\begin{enumerate}[(a)]
  \item \emph{Gaussian linear model for i.i.d.\ pairs.} $y=\alpha^*+x^\top\bstar+\varepsilon$ with
    $\varepsilon\sim\Normal(0,\sigs^2)$, $\sigs^2>0$, independent of $x$; $\alpha^*$, $\bstar$, $\sigs$ and the law of $x$ may
    depend on $n$.
  \item \emph{Sparsity and dimension.} At most $s$ coordinates of $\bstar$ are nonzero, and
    $2\le\praw\le n^A$.
  \item \emph{Sub-Gaussian predictors on sparse directions.} Every predictor has positive variance,
    and $\E\exp(\lambda u^\top\zeta)\le\exp\big(L^2\lambda^2\Var(u^\top\zeta)/2\big)$ for every $\lambda\in\R$ and every
    $u\in\R^{\praw}$ with at most $q_n$ nonzero entries.
  \item \emph{No near-exact collinearity.} $\lmin(\Corr(x_T))\ge c_\psi/n$ for every set $T$ of at
    most $q_n$ predictors.
  \item \emph{Bounded signal.} $\abs{\beta^*_j}\sd(x_j)\le\bar b\,\sigs$ for every $j$.
\end{enumerate}
\end{block}
\end{frame}

\begin{frame}{Assumptions: the reported-model assumptions (1)}
\small
\noindent A \emph{leakage set} is a deterministic $Q$ with $\Sz\subseteq Q\subseteq\set{p}$; its \emph{lattice} is
$\Lat(Q):=\{T:\Sz\subseteq T\subseteq Q\}$. The notation $\Cov(\cdot,\cdot\mid W_U)$ and $\bar\beta_{U,l}$ is that of
Definition~1.3.

\begin{block}{Assumption 5.10 (\asmdisp{V}) \textit{(Visibility: no true predictor masked by $\le15$ others)}}
For every $j\in\Sz$ and every $U\not\ni j$ with $\abs{U}\le15$, $\abs{\Cov(W_j,W^\top\tstar\mid W_U)}\ge a$, for a fixed
$a>0$.
\end{block}
\end{frame}

\begin{frame}{Assumptions: the reported-model assumptions (2)}
\small
\begin{block}{Assumption 5.11 (\asmdisp{reported}) \textit{(Polynomial leakage set with local regressions; A6 Assumption~4)}}
Fix $b\in(0,\tfrac12)$ and $b_H>0$ with $15b_H+b<\tfrac12$.
\begin{enumerate}[(a)]
  \item \emph{Leakage set.} There is a deterministic $Q\supseteq\Sz$ with $\abs{Q}\le n^b$ and, for
    large $n$, $\abs{Q}\le K$ and ($Q=\set{p}$ or $K\ge2\abs{Q}+32\sz+1$), such that: (i) \emph{no leak outside $Q$:}
    there is a deterministic $\varpi_n$ with $\varpi_n\sqrt{n/\log n}\to0$ such that
    $\abs{\Cov(W_l,W_j\mid W_U)}\le\varpi_n$ for all $l\notin Q$, $j\in\Sz\setminus U$ and $\abs{U}\le15$; and (ii)
    \emph{locality:} there is a deterministic $\delta_n\to0$, and every $l\notin Q$ has a deterministic
    $\Hl_l\supseteq\{l\}\cup\Sz$ with $\abs{\Hl_l}\le n^{b_H}$, such that for all $l\notin Q$, all $U\not\ni l$ with
    $\abs{U}\le15$ and $U_0:=U\cap\Hl_l$: $\abs{\Cov(W_l,W_v\mid W_{U_0})}\le\delta_n$ and
    $\abs{\Cov(W_j,W_v\mid W_{U_0})}\le\delta_n$ for all $v\in U\setminus\Hl_l$ and $j\in\Sz\setminus U_0$.
\end{enumerate}
\end{block}
\end{frame}

\begin{frame}{Assumptions: the reported-model assumptions (2, cont.)}
\small
\begin{block}{Assumption 5.11 (\asmdisp{reported}) \textit{(Polynomial leakage set with local regressions; A6 Assumption~4)} (cont.)}
\begin{enumerate}[(a)]\setcounter{enumi}{1}
  \item \emph{Floor on the leakage set.} $\lmin(\Sigma_{n,Q\cup\{l\}})\ge\varphi^2$ for every
    $l\in\set{p}$, with $\varphi$ as in \as{EFS}.
  \item \emph{Local regressions within the leakage set.} There are deterministic sequences
    $d_n=o(\log n)$ and $L_n$, and for every $l\in\set{p}$ a deterministic $N_l\subseteq Q\setminus\{l\}$ with
    $\abs{N_l}\le d_n$, such that $\norm{\bar\beta_{T,l}-\bar\beta_{T\cap N_l,l}}_1\le L_n$ for every $T\in\Lat(Q)$ with
    $l\notin T$, where $L_n\to0$ and $L_n^2\log(\abs{Q}+1)=o(\log n)$.
\end{enumerate}
\end{block}
\end{frame}

\begin{frame}{Assumptions: certified screening}
\small
\begin{block}{Assumption 5.12 (\asmdisp{SCR}) \textit{(Floor on sets of the budget size, and signal bounds; A6 \S5.5)}}
The design is fixed. For a model-size budget $k_n\in\{0,\dots,K-1\}$: $\lmin(\Rhat_S)\ge\kappa$, for a fixed $\kappa>0$,
for every $S$ with $\abs{S}\le k_n+\sz$; $\tmin\ge\tlow>0$ and $\norm{\tstar}_1\le M_1$; $\kappa$, $\tlow$, $M_1$ and
$\sigs^2$ are fixed; and $(k_n+1)\log p+\log n\le c_{\mathrm{SB}}\,n$, where
$c_{\mathrm{SB}}:=\kappa^2\tlow^2/\{200(M_1^2+\sigs^2)\}$.
\end{block}
\end{frame}

\begin{frame}{Assumptions: the witness designs}
\small
\begin{block}{Assumption 5.13 (\asmdisp{WIT}) \textit{(Witness designs; A6 Theorem~6(b))}}
$\sz$ is fixed; the true columns $\Xr_{\Sz}$ and $\tstar$ are deterministic (if they are random, the conditions below hold
conditionally on $(\Xr_{\Sz},\tstar)$ on events of probability tending to one, with the same constants);
$\norm{\mu}^2\le nM_\mu$ and $\norm{\tstar}^2\le nM_\theta$ for fixed $M_\mu$, $M_\theta$; and one of:
\begin{enumerate}[(a)]
  \item \emph{(N-orth)} the columns of $X$ outside $\Sz$ are mutually orthogonal and orthogonal to the
    columns in $\Sz$; or
  \item \emph{(N-gauss)} the raw columns outside $\Sz$ have i.i.d.\ $\Normal(0,1)$ entries, independent
    of the raw true columns $\Xr_{\Sz}$ and of the noise, and $p\le n^A$ for a fixed $A$.
\end{enumerate}
\end{block}
\end{frame}

\begin{frame}{Assumptions of the supplementary results (1)}
\small
\begin{block}{Assumption 5.14 (\asmdisp{BC}) \textit{(Bounded coefficients)}}
$\max_{j}\abs{\theta^*_j}\le\bar\theta$ for a fixed $\bar\theta<\infty$.
\end{block}
\begin{block}{Assumption 5.15 (\asmdisp{EF}) \textit{(Eigenvalue floor on sets of bounded size)}}
$\lmin(\Rhat_S)\ge\varphi_0^2$ for every model $S$ with $\abs{S}\le2(A+1)\sz+2$, for a fixed $\varphi_0\in(0,1]$.
\end{block}
\begin{block}{Assumption 5.16 (\asmdisp{NC}) \textit{(No near-exact collinearity among sets of logarithmic size)}}
$\lmin(G_S)\ge g_0$ for every model $S$ with $\abs{S}\le3(A+1)(\sz+1)\log n$, for a fixed $g_0>0$.
\end{block}
\end{frame}

\begin{frame}{Assumptions of the supplementary results (2)}
\small
\begin{block}{Assumption 5.17 (\asmdisp{BM}) \textit{(Beta-min)}}
Either (a) $\tmin\ge C_\beta\,\sigs\rn$ with a fixed $C_\beta>0$ such that $C_\beta^2\varphi_0^2>c_0(\sz,A)$, the constant
of Definition~5.1; or (b) $\tmin\sqrt{n/\log n}\to\infty$.
\end{block}
\begin{block}{Assumption 5.18 (\asmdisp{margin}) \textit{(Floor on the true model and an identifiability margin; A6 Assumption~A.4)}}
\begin{enumerate}[(a)]
  \item $\lmin(\Rhat_{\Sz})\ge\varphi_0^2$ for a fixed $\varphi_0\in(0,1]$.
  \item Every model $S\in\Om$ that misses a true predictor and has at most $j_0:=\lceil(2A+1)\sz\rceil+1$ predictors
    outside $\Sz$ leaves unexplained signal $\norm{(I-P_S)\mu}^2\ge c_{\mathrm{id}}\,\sigs^2\log n$, for a fixed
    $c_{\mathrm{id}}>c_1(\sz,A)$, the constant of Definition~5.1.
\end{enumerate}
\end{block}
\end{frame}

\begin{frame}{Assumptions of the supplementary results (3)}
\small
\begin{block}{Assumption 5.19 (\asmdisp{omitted}) \textit{(Candidate-only fitting; A6 Assumption~A.6)}}
The method is run on the response $\yr$ and a set $\mathcal C$ of $p$ candidate predictors, with rows $x_{\mathcal C,i}$
collected in $x_{\mathcal C}$; $\Sigma$ and $R$ are the covariance and the correlation matrix of a candidate row.
\begin{enumerate}[(a)]
  \item \emph{Model.} The candidate rows are i.i.d.\ $\Normal_p(m_{\mathcal C},\Sigma)$, and, given them, the responses are
    independent with $\yr_i\mid x_{\mathcal C}\sim\Normal(\alpha'+x_{\mathcal C,i}^\top\beta^{\mathcal C},\sigma^2_{\mathcal C})$,
    $\sigma^2_{\mathcal C}>0$. We write $\Sz^{\mathcal C}:=\supp\beta^{\mathcal C}$ and $s_{\mathcal C}:=\abs{\Sz^{\mathcal C}}$.
  \item \emph{Floor.} $\lmin(R_T)\ge c_{\mathcal C}$ for every $T\subseteq\mathcal C$ with
    $\abs{T}\le3(A+1)(s_{\mathcal C}+1)\log n$, for a fixed $c_{\mathcal C}>0$.
  \item \emph{Signal relative to the residual variance.} $\beta^{\mathcal C\top}\Sigma\beta^{\mathcal C}\le
    M_{\mathcal C}\sigma^2_{\mathcal C}$ for a fixed $M_{\mathcal C}$, and $n\min_{j\in\Sz^{\mathcal C}}\Sigma_{jj}(\beta^{\mathcal
    C}_j)^2/(\sigma^2_{\mathcal C}\log n)\to\infty$.
\end{enumerate}
\end{block}
\end{frame}

\begin{frame}{Assumptions of the supplementary results (4)}
\small
\begin{block}{Assumption 5.20 (\asmdisp{NL}) \textit{(No leakage, as seen by the start law)}}
Every null predictor $l\notin\Sz$ has a deterministic \emph{neighbourhood} $\Hl_l\supseteq\{l\}\cup\Sz$ with
$\abs{\Hl_l}\le e^{\kappa_H\sqrt{\log n}}$ for a fixed $\kappa_H$, such that, for every $U\not\ni l$ with $\abs{U}\le15$ and
$U_0:=U\cap\Hl_l$: (i) \emph{no leak:} $\abs{\Cov(W_l,W_j\mid W_{U_0})}\le C_Ln^{-1/2}$ for every $j\in\Sz\setminus U_0$,
with a fixed $C_L$; and (ii) \emph{locality:} $\abs{\Cov(W_l,W_v\mid W_{U_0})}\le r^\circ_n$ and
$\abs{\Cov(W_j,W_v\mid W_{U_0})}\le r^\circ_n$ for every $v\in U\setminus\Hl_l$ and $j\in\Sz\setminus U_0$, where
$r^\circ_n(n\log n)^{1/4}\to0$.
\end{block}
\end{frame}

\begin{frame}{Assumptions of the supplementary results (5)}
\small
\noindent For a set $U$ and a predictor $l\notin U$: the \emph{signal gain} is $\delta^2_{U,l}:=\norm{(P_{U\cup\{l\}}-P_U)\mu}^2$;
the \emph{residual scale} is $\tau^2_U:=\sigs^2+\norm{(I-P_U)\mu}^2/(n-1)$; and $\hat\beta_{U,l}:=(X_U^\top X_U)^{-1}X_U^\top
x_l$, padded with zeros.

\begin{block}{Assumption 5.21 (\asmdisp{reported-fixed}) \textit{(Fixed-design visibility and leakage; A6 Assumption~A.13)}}
\begin{enumerate}[(a)]
  \item \emph{Floors.} $\lmin(\Rhat_T)\ge\varphi_0^2$ for every $T$ with $\abs{T}\le16$ and for every
    $T=Q\cup\{l\}$, $l\in\set{p}$, with $Q$ as in~(c).
  \item \emph{Visibility.} There is a deterministic $\omega_n\to\infty$ such that
    $\delta^2_{U,j}\ge\omega_n\log n$ for every $j\in\Sz$ and every $U\not\ni j$ with $\abs{U}\le15$.
\end{enumerate}
\end{block}
\end{frame}

\begin{frame}{Assumptions of the supplementary results (5, cont.)}
\small
\begin{block}{Assumption 5.21 (\asmdisp{reported-fixed}) \textit{(Fixed-design visibility and leakage; A6 Assumption~A.13)} (cont.)}
\begin{enumerate}[(a)]\setcounter{enumi}{2}
  \item \emph{Bounded leakage.} There are a deterministic $Q\supseteq\Sz$ (with $\abs{Q}\le K$ and
    ($Q=\set{p}$ or $K\ge2\abs{Q}+32\sz+1$) for large $n$) and a fixed $\eta\in[0,\eta_A)$,
    $\eta_A:=(\sqrt{2A+1}-\sqrt{2A})^2$, such that $\delta^2_{U,l}\le\eta\,\tau^2_U\log n$ for every $l\notin Q$ and every
    $U\not\ni l$ with $\abs{U}\le15$.
  \item \emph{Size of the leakage set.} One of: (d1) $\abs{Q\setminus\Sz}\le(1-\delta_Q)\log n$ for a
    fixed $\delta_Q\in(0,1)$; or (d2) $\abs{Q}\le n^b$ with $b<\tfrac12-\tfrac\eta2-\sqrt{2A\eta}$,
    $\theta^{*2}_{\min}n/(\abs{Q}+\log n)\to\infty$, and local regressions: for every $l$ a deterministic set
    $N_l\subseteq Q\setminus\{l\}$ with $\abs{N_l}\le h\log_2n$, for a fixed $h$ with $h+b<\tfrac12$, and a deterministic
    $L_n\to0$ with $L_n^2\log(\abs{Q}+1)=o(\log n)$, such that $\norm{\hat\beta_{T,l}-\hat\beta_{T\cap N_l,l}}_1\le L_n$
    for every $T\in\Lat(Q)$ with $\abs{T}<K$ and $l\notin T$.
\end{enumerate}
\end{block}
\end{frame}

\begin{frame}{Assumptions of the supplementary results (6): explicit constants}
\small
\begin{block}{Definition 5.1 (Explicit constants)}
With $j_0:=\lceil(2A+1)\sz\rceil+1$ and $C_x:=Aj_0+\sz+\log(j_0+1)+2$:
$c_0(\sz,A):=\max\{4C_*,\,32[\sz(A+1)+2]\}$ with $C_*:=2\sz+3[Aj_0+\sz+\log(j_0+1)+1]=3C_x+2\sz-3$; and, with
$C_A:=8+3A$,
\[
  c_1(\sz,A):=\Big(\sqrt{2C_x}+\sqrt{2C_x+C_Aj_0+(2/\log2)\{\sz(A+1)+2\}}\Big)^2 .
\]
\end{block}
\end{frame}

\begin{frame}{Conventions (1)}
\small
\noindent The conventions fix what ``the method'' means and idealise the computer; they are not statistical assumptions.

\begin{block}{Convention 5.22 (\asmdisp{C1}) \textit{(The method)}}
Exact arithmetic, the default configuration, valid inputs and seed. A run that meets it is \emph{admissible}; a sample
size $n$ is admissible if the method accepts some data set of that size as valid input. A statement that holds for every
valid seed and every outcome of the calibration says so.
\end{block}
\end{frame}

\begin{frame}{Conventions (2): delivery}
\small
\begin{block}{Convention 5.23 (\asmdisp{C2}) \textit{(Delivery)}}
\begin{enumerate}[(i)]
  \item \emph{Idealised machine:} exact arithmetic at unit cost (including $\exp$ and $\log$); no integer, index,
    random-stream identifier or allocation overflow; wall time at most $c_{\mathrm t}\cdot(\text{operation count})+c_{\mathrm
    t0}$ and peak memory at most $c_{\mathrm m}\cdot(\text{words held})+c_{\mathrm m0}$ for fixed constants $(c_{\mathrm
    t},c_{\mathrm t0},c_{\mathrm m},c_{\mathrm m0})$; every allocation within the declared memory succeeds; the build and its
    dependencies work; no external interruption or worker crash.
\end{enumerate}
\end{block}
\end{frame}

\begin{frame}{Conventions (2, cont.)}
\small
\begin{block}{Convention 5.23 (\asmdisp{C2}) \textit{(Delivery)} (cont.)}
\begin{enumerate}[(i)]\setcounter{enumi}{1}
  \item \emph{Covered configuration:} the default core of the software, with the local search from the start mode (cap
    $\ell_n+\abs{\tSz}$) and the rule that resolves ties in the start mode, implemented with
    $O\big(p\sum_k(2^{\abs{C_k}}+(K+1)^2)\big)$ operations and $O\big(\sum_k2^{\abs{C_k}}+J(K+1)\big)$ storage; the local
    search from the start mode starting at the smallest, lexicographically first maximiser of the start law; the default
    list of summaries only; certificate computations with declared polynomial operation and memory budgets, included in
    the allocation; serial execution or a fixed number of concurrent workers (memory envelope multiplied by that number);
    a fixed number of islands; the fallback temperature sequence switched on (the default).
\end{enumerate}
\end{block}
\end{frame}

\begin{frame}{Conventions (2, cont.)}
\small
\begin{block}{Convention 5.23 (\asmdisp{C2}) \textit{(Delivery)} (cont.)}
\begin{enumerate}[(i)]\setcounter{enumi}{2}
  \item \emph{Random numbers and stops:} ideal independent uniform random numbers; the method's stop-and-return rule and
    stop-only resource guard as coded (a stopped run returns no estimates).
  \item \emph{Inputs and allocation:} a valid non-degenerate input (finite data, a response of length $n\ge3$ that is not
    constant, at least one non-constant predictor column), a valid seed (an integer in $[1,2^{31}-1]$), and a declared memory and time
    allocation at least the method's explicit polynomial work and memory envelopes in $(n,\praw)$, including the memory
    checks made before the run and the bounds on the records and time of the local searches.
\end{enumerate}
\end{block}
\end{frame}

\begin{frame}{Conventions (3)}
\small
\begin{block}{Convention 5.24 (\asmdisp{C2p}) \textit{(Delivery, with the integer limits lifted)}}
Convention~\as{C2}, with in addition the two fixed integer limits of Table~1 (the ceiling $10^6$ on $\ell_n$
and the bound $2^{31}-1$ on $\ell_n+\abs{\tSz}$) idealised away; every other validator and cap (including the cap on the
enumeration used to find the start mode, and its fallback) is kept as coded.
\end{block}
\end{frame}

\begin{frame}{Conventions (4)}
\small
\begin{block}{Convention 5.25 (\asmdisp{C3}) \textit{(Sampler)}}
\begin{enumerate}[(i)]
  \item The computer does exact arithmetic.
  \item Within each random stream, the random numbers are fresh, independent, ideal uniforms; the final run's stream for
    each island is independent of the data and of the calibration's streams.
  \item The coded start sampler returns exact draws from $\pi_0$, and the reported start normaliser equals the exact
    normalising constant of the unnormalised weights of $\pi_0$.
  \item The coded Markov moves realise their stated kernels, applied in a uniformly random order that is independent of
    the state.
  \item The outcome of the calibration is a function of the data and of the calibration's random streams, and each island's
    record is a measurable function of that outcome and of the island's stream.
  \item Statements about estimates refer to the potential (uninterrupted) run, which is identified with the reported
    values on completed runs.
\end{enumerate}
\end{block}
\end{frame}

\begin{frame}{Which statements use which assumption}
\footnotesize
{\footnotesize Table 2. Assumptions and conventions against the statements that use them, read from the hypotheses of each statement.
Theorems~1--6 are in Section~6.1; S-numbers are in Section~7.}
\smallskip

\begin{tabular}{@{}>{\raggedright\arraybackslash}p{0.27\linewidth}>{\raggedright\arraybackslash}p{0.68\linewidth}@{}}
\toprule
Assumption & Used by \\
\midrule
Assumption~\as{GL} & Thm.~1, 4, 6(b); S2, S3, S4, S11 \\
Assumption~\as{DIM} & Thm.~1, 3; S2, S3, S4, S5, S6, S10, S11, S12, S14, S15 \\
Assumption~\as{SP} & Thm.~1; S2, S3, S4, S6, S11 \\
Assumption~\as{GR} & Thm.~3; S10, S12, S14, S15 \\
Assumption~\as{SGR} & Thm.~1 \\
Assumption~\as{FSb} & Thm.~1 \\
Assumption~\as{FSt} & Thm.~3; S10, S12, S14, S15 \\
Assumption~\as{EFS} & Thm.~1, 3; S14, S15 \\
Assumption~\as{dist} & Thm.~2; S7, S8, S9; planned P3 \\
\bottomrule
\end{tabular}
\end{frame}

\begin{frame}{Which statements use which assumption (cont.)}
\footnotesize
\begin{tabular}{@{}>{\raggedright\arraybackslash}p{0.27\linewidth}>{\raggedright\arraybackslash}p{0.68\linewidth}@{}}
\toprule
Assumption & Used by \\
\midrule
Assumption~\as{V} & Thm.~3; S10, S14, S15 \\
Assumption~\as{reported} & Thm.~3; S14; planned P2(i), P3, P4 \\
Assumption~\as{SCR} & Thm.~4 \\
Assumption~\as{WIT} & Thm.~6(b) \\
Assumption~\as{BC} & S2, S3, S4, S11 \\
Assumption~\as{EF} & S2, S3, S11 \\
Assumption~\as{NC} & S2, S3, S4, S11 \\
Assumption~\as{BM} & S2, S11 \\
Assumption~\as{margin} & S4 \\
\bottomrule
\end{tabular}
\end{frame}

\begin{frame}{Which statements use which assumption (cont.)}
\footnotesize
\begin{tabular}{@{}>{\raggedright\arraybackslash}p{0.27\linewidth}>{\raggedright\arraybackslash}p{0.68\linewidth}@{}}
\toprule
Assumption & Used by \\
\midrule
Assumption~\as{omitted} & S5 \\
Assumption~\as{NL} & S10, S15 \\
Assumption~\as{reported-fixed} & S11 \\
Convention~\as{C1} & Thm.~3, 4; S11, S12, S14, S15 \\
Convention~\as{C2} & Thm.~3; S11, S12 \\
Convention~\as{C2p} & Thm.~3; S11, S12, S14 \\
Convention~\as{C3} & Thm.~5; S15, S16, S17 \\
No assumption & Thm.~6(a); S1; S13 (gives \as{V}) \\
\bottomrule
\end{tabular}
\end{frame}

\section{The proved theorems}
\begin{frame}
\centering\Large\textbf{The proved theorems}
\end{frame}

\begin{frame}{Theorem 1}
\small
\begin{block}{Theorem 1 (Posterior concentration, random sub-Gaussian design)}
Assume \as{GL}, \as{DIM}, \as{SP}, \as{SGR}, \as{FSb} and \as{EFS}, where in \as{DIM}, \as{SGR} and \as{EFS} the
dimension is the number $\praw$ of raw predictor columns (before the method removes constant columns) and the sets are
sets of raw columns; the prior, the cap $K$ and the posterior use the retained columns. Columns keep their raw labels, and $\Sz:=\{j\in\set{\praw}:\beta^*_j\neq0\}$ with $\bstar$ the coefficients of
the raw design; on the event that the method removes a column, set $\PIP_j:=0$ for a removed column $j$ and
$\pis{\Sz}:=0$ if $\Sz\notin\Om$, and, if no column is retained, take the posterior to be a point mass at $\emptyset$.
Then, under the joint law of the design and the noise,
\begin{gather}
  1-\pis{\Sz}=\OP\Big(\sqrt{\frac{\log n}{n}}\Big),\qquad \P\big(\HPM=\MPM=\Sz\big)\to1,\tag{12}\\
  \max_{j\in\set{\praw}}\big|\PIP_j-\ind{j\in\Sz}\big|=\OP\Big(\sqrt{\frac{\log n}{n}}\Big).\nonumber
\end{gather}
\end{block}
\medskip\noindent\textbf{Status: PROVED}\quad{\footnotesize Registry: \rid{W3-D3:Thm2SG}, \rid{W3-D3:B3w}, \rid{W3-D3:ExSG}: \rid{PROVED-CHECKED}. Gaussian rows ($L=1$): \rid{W3-POST:Thm-1}, \rid{W2-APP:POST}, \rid{W3-READER:R1}: \rid{PROVED-CHECKED}.}
\end{frame}

\begin{frame}{Theorem 2}
\small
\begin{block}{Theorem 2 (Distribution-level consistency)}
There are constants $c_{\mathrm H}>0$, $M_0$ and $n_0$, depending only on $(s,A,L,c_\psi,\bar b)$, such that the following
holds. Fix a level $M\ge M_0$ and assume Assumption~\as{dist}, with the window $q_n$ of this $M$. On a data set on which fewer
than two predictors remain or the response is constant, the posterior is not defined, and the posterior probabilities
below are read as $1$. Then for every $n\ge n_0$ with $\sqrt n/\log n\ge M$,
\begin{equation}\tag{13}
  \P^*\Big[\,\Pi\Big(d_H(f,f^*)>\sqrt{M\log n/n}\ \Big|\ \Dn\Big)>n^{-c_{\mathrm H} M}\Big]\ \le\ n^{-c_{\mathrm H} M}.
\end{equation}
\end{block}
\medskip\noindent\textbf{Status: PROVED}\quad{\footnotesize Registry: \rid{W3-DIST:Thm-A}: \rid{PROVED-CHECKED}.}
\end{frame}

\begin{frame}{Theorem 3}
\small
\begin{block}{Theorem 3 (The reported model equals the true model)}
Let $\Sig$ be positive definite and assume \as{GR}, \as{DIM}, \as{FSt}, \as{V}, \as{EFS}, Assumption~\as{reported}
and Convention~\as{C1}. Then, under the joint law of the design and the noise, there are data events $\mathcal G_n$ with
$\P(\mathcal G_n)\to1$ on which \emph{every} admissible run that produces output reports $\Shat=\Sz$, whatever its seed and
the outcome of its internal calibration. In particular,
\begin{equation}\tag{14}
  \P(\text{the run produces output and }\Shat\neq\Sz)\to0 .
\end{equation}
If moreover Convention~\as{C2} holds, then $\P(\text{the run produces output})=1$ for every admissible $(n,\praw)\in\Dom$
(Definition~4.11). If Convention~\as{C2p} holds, then
\begin{equation}\tag{15}
  \P(\text{the run produces output})=1\ \ \text{for every admissible }n\ge3,\qquad \P(\Shat=\Sz)\to1 .
\end{equation}
\end{block}
\medskip\noindent\textbf{Status: PROVED}\quad{\footnotesize Registry: \rid{W3-RM2-R:RM2-R-plus}, \rid{W3-RM2-R:RM0-plus}, \rid{W3-RM2-R:DLV-0}: \rid{PROVED-CHECKED}. Code conformance \rid{C-DLV:L2-S3}, \rid{C-DLV:P1-Inventory}: \rid{CODE-CHECKED (second seat, code reading; holds at pin 4676af6 after amendments A-1--A-7)}.}
\end{frame}

\begin{frame}{Theorem 4}
\small
\begin{block}{Theorem 4 (Certified screening)}
Let the design be fixed and assume \as{GL}, Assumption~\as{SCR} with model-size budget $k_n$, and Convention~\as{C1};
\as{SP} is not assumed. Then
\[
  \P\big(\text{the run produces output},\ \Shat\text{ is certified},\ \abs{\Shat}\le k_n,\ \Sz\not\subseteq\Shat\big)\to0 .
\]
\end{block}
\medskip\noindent\textbf{Status: PROVED}\quad{\footnotesize Registry: \rid{W3-SCR:A1}, \rid{W3-SCR:A-inf}, \rid{W3-SCR:SCR-A}: \rid{PROVED-CHECKED}.}
\end{frame}

\begin{frame}{Proposition 5}
\small
\begin{block}{Proposition 5 (Each island's evidence estimate is unbiased)}
Assume Convention~\as{C3}, and let $\mathcal C$ be the calibration record. For every fixed data set (probability over the
final run's random streams), almost surely on $\Acal$, each island's potential evidence estimate satisfies
$\E[\Zhat_r\mid\mathcal C]=Z$, $r=1,\dots,4$, and hence the pooled estimate~\textup{(11)} satisfies
$\E[\Zhat\mid\mathcal C]=Z$. Under the joint law of the data and the method's random numbers, the same identities hold
conditionally on the data and $\mathcal C$, and $\E[\ind{\Acal}\Zhat]=\E[\ind{\Acal}Z]$. No condition on the design is
made.
\end{block}
\medskip\noindent\textbf{Status: PROVED}\quad{\footnotesize Registry: \rid{W1-N1:Thm-N1}, \rid{W3-READER:PoolZ}, \rid{W3-AUDIT-AX:FR}: \rid{PROVED-CHECKED}; line checks \rid{W3-AUDIT-AW:ISL}, \rid{W3-AUDIT-AW:FRZ}.}
\end{frame}

\begin{frame}{Theorem 6}
\small
\begin{block}{Theorem 6 (The rate cannot be improved)}
\textup{(a)} \emph{Every design.} For every $n\ge3$, every standardised design with $p\ge2$, every response and every
model $S$ with $\abs{S}\le\min(K-1,p/2)$,
\begin{equation}\tag{16}
  \pis{S}\ \le\ \Big(1+\frac{\abs{S}+1}{3\sqrt n}\Big)^{-1}.
\end{equation}
So no model of that size, the true one included, can carry posterior mass closer to one than about
$(\abs{S}+1)/(3\sqrt n)$. For true models with $\sz\le\min(K-1,p/2)$, no design gives a rate faster than
$n^{-1/2}$, and a posterior miss of order $\rn$ is possible only if $\sz=O(\sqrt{\log n})$.
\end{block}
\end{frame}

\begin{frame}{Theorem 6 (cont.)}
\small
\begin{block}{Theorem 6 (The rate cannot be improved) (cont.)}
\textup{(b)} \emph{Witness designs.} Assume \as{GL} and \as{WIT}. Then, for every $p>\sz$, bounded or growing, with
$\pi=3.14\ldots$ in this display,
\begin{equation}\tag{17}
  \E\big[1-\pis{\Sz}\big]\ \ge\ (1-o(1))\,(\sz+1)\,\frac{2(p-\sz)}{2p-1-\sz}\,
  \sqrt{\frac{\log\!\big(p\sqrt n/(\sz+1)\big)}{\pi n}},
\end{equation}
and, if moreover $p\to\infty$, $1-\pis{\Sz}\ge(1-\oP(1))\,(\sz+1)\sqrt{\log p/(\pi n)}$.
\end{block}
\medskip\noindent\textbf{Status: PROVED}\quad{\footnotesize Registry: (a)~\rid{W3-PSHARP:B}; (b)~\rid{W3-PSHARP:A} (with Corollary~A1) and \rid{W3-PSHARP-R:AE} (with Corollary~A1-E):
\rid{PROVED-CHECKED}.}
\end{frame}

\section{The theorems to prove (PLANNED)}
\begin{frame}
\centering\Large\textbf{The theorems to prove (PLANNED)}
\end{frame}

\begin{frame}{Definition 6.1 (used by P1)}
\small
\begin{block}{Definition 6.1 (Single-flip chains, mixing time)}
A \emph{single-flip chain} is a Markov chain on $\Om$ that leaves $\pis{\cdot}$ invariant and whose every transition
adds or removes at most one predictor. For a Markov kernel $P$ on $\Om$ with invariant law $\pis{\cdot}$, the
\emph{mixing time} is $t_{\mathrm{mix}}(P):=\min\{t\ge0:\max_{S\in\Om}\norm{P^t(S,\cdot)-\pis{\cdot}}_{\mathrm{TV}}\le1/4\}$.
The \emph{weighted tempered Gibbs sampler} is the sampler of Zanella and Roberts (2019) targeting $\pis{\cdot}$.
\end{block}
\end{frame}

\begin{frame}{Planned Theorem P1 --- PLANNED}
\small
\begin{alertblock}{Planned Theorem P1 (The trap theorem; candidate)}
There is a family $\mathfrak F$ of design sequences \textup{[open choice~1]} whose posterior has several well-separated
modes, such that for every sequence in $\mathfrak F$:
\begin{enumerate}[(i)]
  \item with probability tending to one, $t_{\mathrm{mix}}(P)\ge\tau_n$ for every single-flip chain $P$
    \textup{[open choice~3]} and for the weighted tempered Gibbs sampler, where $\tau_n\to\infty$ at a rate
    \textup{[open choice~4]} that grows quickly with the problem size; and
  \item the hypotheses of Theorem~$k$ hold, for one $k\in\{1,3,4\}$ \textup{[open choice~2]}, so that its conclusion holds
    on $\mathfrak F$.
\end{enumerate}
\end{alertblock}
\medskip\noindent\textbf{Status: PLANNED}\quad{\footnotesize No statement exists in any project file. Owner: Theory Department (on hold). Order: after the $p=10^6$ mission,
before Planned Theorem~P2.}
\end{frame}

\begin{frame}{Planned Theorem P1 --- open choice}
\small
\begin{block}{Open choice P1 (not decided; see Appendix B of the specification)}
(1)~The design family $\mathfrak F$: no statement, draft or design family exists. (2)~Which result is
``our consistency on the same design'': Theorem~1, 3 or~4. If it is
Theorem~3, a leakage set $Q$ satisfying Assumption~\as{reported} must be exhibited for $\mathfrak F$
(the leakage-set choice). (3)~The class of chains: every single-flip chain, or one named kernel (for example a
Metropolis--Hastings add/remove chain). (4)~The growth rate of $\tau_n$ (for example exponential in $n$ or in $\sz$).
\end{block}
\end{frame}

\begin{frame}{Planned Theorem P2 --- PLANNED}
\small
\begin{alertblock}{Planned Theorem P2 (Coverage at $p=10^7$; candidate)}
Let $A_7:=\log(10^7)/\log200\ (\approx3.04)$, and let $\mathfrak D_{\mathrm{AR}}$ be the autoregressive showcase design,
continued as a sequence in $n$ with $p_n=\lceil n^{A_7}\rceil$, a known true model and its structure fixed
\textup{[open choice~2]}.
\begin{enumerate}[(i)]
  \item \emph{(Reading ``theory coverage''.)} $\mathfrak D_{\mathrm{AR}}$ satisfies the hypotheses of Theorems
    \textup{[open choice~3]} with $A=A_7$; in particular, for Theorem~3, there is a leakage set $Q$
    satisfying Assumption~\as{reported} \textup{[open choice~4]}.
  \item \emph{(Reading ``claim coverage''.)} Of the statements of Sections~6.1 and~7, those
    whose conclusion holds for every seed on a data event (Theorem~3, Theorem~S11,
    Corollary~S14, Lemma~S1, Theorem~6(a)) are the ones a single fit with one
    seed can support.
\end{enumerate}
\end{alertblock}
\medskip\noindent\textbf{Status: PLANNED}\quad{\footnotesize The engine design is paused; the theory analysis is not started. Owner: Theory (analysis), Engineering (engine).
Order: after Planned Theorem~P1.}
\end{frame}

\begin{frame}{Planned Theorem P2 --- open choice}
\small
\begin{block}{Open choice P2 (not decided; see Appendix B of the specification)}
(1)~What ``coverage at $p=10^7$'' means: reading~(i), whether the hypotheses of the theorems can hold at
$n=200$, $p=10^7$; or reading~(ii), which statements a single fit with one seed can support (reading~(ii) may not be a
theorem). No project document defines ``coverage''. (2)~The showcase design: the owner fixed one simulated AR-strong data set with $n=200$, a known truth, one fit and one
seed (\code{HQ_DECISIONS.md}); the continuation as a sequence in $n$ is not fixed. (3)~Which theorems reading~(i) covers. (4)~The leakage
set $Q$ for $\mathfrak D_{\mathrm{AR}}$. Also open: the fit runs on an engine derived from the $p=10^6$ engine
(Section~4.4), which no theorem covers.
\end{block}
\end{frame}

\begin{frame}{Planned Theorem P3 --- PLANNED}
\small
\begin{alertblock}{Planned Theorem P3 (The raw-scale search bridge; candidate)}
Assume Assumption~\as{dist} with $s\ge1$ and the window $q_n$ of a level $M\ge M_0$, and assume raw-scale counterparts of
the hypotheses of Theorem~3 \textup{[open choice~1]}, so that the true coefficients $\bstar$ are fixed on
the raw scale and the standardised coefficients $\tstar=D\bstar$ are random. Then there are data events of
$\P^*$-probability tending to one on which every admissible completed run, whatever its seed and calibration outcome,
returns a list $\mathcal T$ with $\pi(\mathcal T\mid\Dn)\ge n^{-c_{\mathrm H}M/2}$. Consequently the conclusion of
Remark~S9 holds for every such run on events of $\P^*$-probability tending to one
\textup{[open choice~2]}.
\end{alertblock}
\medskip\noindent\textbf{Status: PLANNED}\quad{\footnotesize Named open item (\code{QUEUE.md}, ``Named open items'', D121; \code{READER_VERSIONS.md}, row B6.7). Not a ticket;
no statement in any project file. Owner: Theory (on hold).}
\end{frame}

\begin{frame}{Planned Theorem P3 --- open choice}
\small
\begin{block}{Open choice P3 (not decided; see Appendix B of the specification)}
(1)~The raw-scale hypotheses: Theorem~3 and its masters take $\tstar$ deterministic on the
standardised scale, while Theorem~2 takes i.i.d.\ pairs with fixed raw $\bstar$; for $\sz\ge1$ no design
satisfies both (erratum E3 of \code{W3-D5P}), so the population conditions (\as{FSt}, \as{V}, Assumption~\as{reported})
must be restated for fixed $\bstar$. (2)~The target: the list-mass premise only (as stated), or the full conclusion of
Theorem~3 on the raw scale. Today the premise is known only pathwise, on the part of the data event of
Corollary~S14 where $\pi(\Sz\mid y)\ge n^{-c_{\mathrm H}M/2}$.
\end{block}
\end{frame}

\begin{frame}{Planned Theorem P4 --- PLANNED}
\small
\begin{alertblock}{Planned Theorem P4 (The reported model under an aggregated leak; candidate)}
The conclusions of Theorem~3 hold when, in Assumption~\as{reported}, the per-predictor bound
(a)(i) is replaced by a condition on the leaks aggregated over the predictors outside $Q$
\textup{[open choice~1]}, the other hypotheses of Theorem~3 being kept.
\end{alertblock}
\medskip\noindent\textbf{Status: PLANNED}\quad{\footnotesize Named open item (\code{QUEUE.md}, ``Named open items''; \code{W3-RM2-R/STATEMENT.md}, ``Not targeted''). No
statement in any project file. Owner: Theory (on hold).}
\begin{block}{Open choice P4 (not decided; see Appendix B of the specification)}
(1)~The aggregated condition: no form is written in any project file (\code{W3-RM2-R/STATEMENT.md} and
\code{PLAN.md} name ``the aggregated Form~R leak'' only). (2)~Whether it is stated for the random-design form
(Theorem~3) only, or also for the fixed-design form (Theorem~S11).
\end{block}
\end{frame}


% P5 and Target N (manager). Source: Theory_work/Dept/assess/MIS_KL/ASSESS_fable.md §2 (setting (P), Target A, Target N);
% open-choice block: ASSESS_astra.md §2, §3 (intercept), §5. Statements verbatim; notation of the ASSESS files.
\begin{frame}{P5 (new) --- Misspecification under a KL projection, fixed $p$: setting}
\begin{alertblock}{Setting (P)}
Rows $(x_i,y_i)$ i.i.d.\ from a fixed law $P^*$ on $\mathbb R^p\times\mathbb R$, $\Sigma=\mathrm{Var}(x)\succ0$,
$\beta^0=\Sigma^{-1}\mathrm{Cov}(x,y)$, $\alpha^0=Ey-Ex^\top\beta^0$, $r=y-\alpha^0-x^\top\beta^0$, $v=Er^2>0$,
$S_0=\operatorname{supp}\beta^0$, $f_0=\nu_x(dx)\,\phi(y;\alpha^0+x^\top\beta^0,v)$.
\end{alertblock}
\end{frame}

\begin{frame}{P5 (new) --- Misspecification under a KL projection, fixed $p$}
\small
\begin{alertblock}{Planned Theorem P5 (Target A; candidate)}
Under (P) with $p$ fixed, $E[\|x\|^2r^2]<\infty$ and $Er^4<\infty$, for the coded posterior:
\[
1-\pi(S_0\mid D_n)=O_P(n^{-1/2}),\quad P(\mathrm{HPM}=\mathrm{MPM}=S_0)\to1,\quad
\max_j|\mathrm{PIP}_j-1\{j\in S_0\}|=O_P(n^{-1/2}),
\]
\[
\Pi\big(d_H(f_{S,\theta,v'},f_0)>M_nn^{-1/2}\ \text{or}\ \mathrm{KL}(f_0\,\|\,f_{S,\theta,v'})>M_n/n\ \big|\ D_n\big)\to0
\ \text{in probability, for every }M_n\to\infty,
\]
and $f_0$ is the Kullback--Leibler projection of $P^*$ whenever that divergence is finite for some member of the family.
No Gaussianity, no linear conditional mean, no homoscedasticity.
\end{alertblock}
\textbf{Status: PLANNED} \quad{\footnotesize Source: \texttt{Theory\_work/Dept/assess/MIS\_KL/ASSESS\_fable.md} \S2.}
\end{frame}

\begin{frame}{P5 --- open choice (Astra's lock conditions)}
\small
\begin{block}{Open choice P5 (not decided)}
\begin{itemize}
  \item \textbf{Exact selection or consistency.} ``Exactness is plausible and worth targeting; inclusion plus Hellinger
    consistency is a fallback'': $\Pi\{H(f,f_0)>\epsilon\mid D_n\}\to_{P_*}0$ $(\epsilon>0)$, $\pi(S_0\mid D_n)\to_{P_*}1$.
  \item \textbf{Literal KL.} ``Assume a unique projection, a positive-definite covariance of the retained predictors, finite
    joint fourth moments, and finite $\mathrm{KL}(P_*\Vert f_0)$ for the literal KL interpretation. Otherwise describe the
    target as the Gaussian expected-log-score projection.''
  \item \textbf{Intercept.} ``With raw coefficients and the plug-in intercept of \texttt{W3-DIST/STATEMENT.md} \S0.''
    ``A common random intercept cancels for model selection, but can destroy iid rows and a deterministic population
    density target; fix it initially.''
  \item \textbf{Example.} ``Require an example outside (OV$^{c}$), and an explicit audit of standardisation and the
    variance prior.''
\end{itemize}
\end{block}
{\footnotesize Source: \texttt{Theory\_work/Dept/assess/MIS\_KL/ASSESS\_astra.md} \S2, \S3, \S5.}
\end{frame}

\begin{frame}{Target N --- the heteroscedastic boundary at growing $p$}
\begin{alertblock}{Target N (candidate)}
Fixed design, $y^r=\alpha1+x\beta+\varepsilon$, $\varepsilon\sim N_n(0,\mathrm{diag}(v_i))$, all design hypotheses of
MIS-R (i) in force. For a null column put
$\rho_j=\sum_i\tilde x_{ij}^2v_i\big/\big(\bar v\sum_i\tilde x_{ij}^2\big)$, with $\tilde x_j$ the part of $x_j$ orthogonal
to $x_{S_0}$ and $\bar v$ the mean variance. If $p=n^{A'}$ and $p^{1-o(1)}$ mutually orthogonal null columns have
$\rho_j\ge\rho>1+1/(2A')$, then $\pi(S_0\mid x,y^r)\to0$ in probability.
\end{alertblock}
\textbf{Status: PLANNED, conditional: kept only if a valid design exists.}
\quad{\footnotesize Source: \texttt{ASSESS\_fable.md} \S2.}
\end{frame}


\section{The supplementary statements}
\begin{frame}
\centering\Large\textbf{The supplementary statements}
\end{frame}

\begin{frame}{Supplementary: the posterior (S1, S3)}
\footnotesize
\begin{block}{Lemma S1 (The posterior miss controls the summaries)}
Set $\pis{\Sz}:=0$ if $\Sz\notin\Om$. For every data set, $\max_{j\in\set{p}}\abs{\PIP_j-\ind{j\in\Sz}}\le1-\pis{\Sz}$. If $\pis{\Sz}>1/2$, then
$\HPM=\MPM=\Sz$.
\end{block}
{\footnotesize \textbf{PROVED}\quad Registry: \rid{C-PACK:SL0}: \rid{PROVED-CHECKED}.}

\begin{block}{Remark S3 (Without beta-min)}
Under the other assumptions of Theorem~S2, the posterior mass of models that strictly contain $\Sz$ is
$\OP(\rn)$, and every true predictor with $\theta_j^{*2}\ge c\,\sigs^2\log n/(n\varphi_0^2)$, $c>c_0(\sz,A)$, has
$1-\PIP_j=\OP(\rn)$.
\end{block}
{\footnotesize \textbf{PROVED}\quad Registry: \rid{W3-POST:Thm-1} (item 5): \rid{PROVED-CHECKED}.}

\end{frame}

\begin{frame}{Supplementary: the posterior (S2)}
\footnotesize
\begin{block}{Theorem S2 (Posterior concentration, fixed design)}
Let the design be fixed and assume \as{GL}, \as{DIM}, \as{SP}, \as{BC}, \as{EF}, \as{NC} and \as{BM}. Then, along every
sequence that satisfies these assumptions,
\begin{equation}\tag{18}
  \max_{j\in\set{p}}\big|\PIP_j-\ind{j\in\Sz}\big|\ \le\ 1-\pis{\Sz}\ =\ \OP\Big(\sqrt{\frac{\log n}{n}}\Big),\qquad
  \P\big(\HPM=\MPM=\Sz\big)\to1 .
\end{equation}
\end{block}
\end{frame}

\begin{frame}{Supplementary: the posterior (S2) (cont.)}
\footnotesize
\begin{block}{Theorem S2 (Posterior concentration, fixed design) (cont.)}
Moreover, there are $n_0$ and $C_{\mathrm E}$, depending only on $(\sz,A,\sigs^2,\bar\theta,\varphi_0,g_0,C_\beta)$, such that for
every $n\ge n_0$ and every design and $\tstar$ that satisfy the assumptions at that $n$,
\begin{equation}\tag{19}
  \E\big[1-\pis{\Sz}\big]\ \le\ C_{\mathrm E}\sqrt{\frac{\log n}{n}},\qquad \P\big(\HPM=\MPM=\Sz\big)\ \ge\ 1-4C_{\mathrm E}\sqrt{\frac{\log n}{n}} .
\end{equation}
The expectation is over the noise, with the design fixed. Under form (b) of \as{BM}, $n_0$ also depends on how fast
$\tmin\sqrt{n/\log n}$ diverges.
\end{block}
{\footnotesize \textbf{PROVED}\quad Registry: In probability: \rid{W3-POST:Thm-1}: \rid{PROVED-CHECKED}; \rid{SOP:TH09}: \rid{PROVED-CHECKED (T03$\sharp$ only: the conjunction of W3-T03ACC:T03S.1--T03S.4, clause map
CHECKPOINT\_4.md \S2; \dots). Not upgraded (their statuses stay as at theory\_main.tsv:68): T02(i$\sharp$) and
T02(ii$\sharp$) PROVED, SF/XF VERIFIED; T01--T03 PROVED-LOCKED (wave 3; external review queued), SAME-FAMILY-ONLY}. Expectation: \rid{W3-POST2:C}: \rid{PROVED-CHECKED, conditional on W2-APP:POST}; with \rid{W3-POST:Prop-LB}, \rid{W3-POST:Lemma-1a\textquotesingle}: \rid{PROVED-CHECKED}.}

\end{frame}

\begin{frame}{Supplementary: the posterior (S4)}
\footnotesize
\begin{block}{Theorem S4 (Posterior concentration under an identifiability margin)}
Let the design be fixed and assume \as{GL}, \as{DIM}, \as{SP}, \as{BC}, \as{NC} and Assumption~\as{margin}. Then the
conclusions~\textup{(18)}--\textup{(19)} of Theorem~S2 hold, the expectation bound with
constants depending only on $(\sz,A,\sigs^2,\bar\theta,\varphi_0,g_0,c_{\mathrm{id}})$. Under \as{EF} and \as{BM}(a), every model of Assumption~\as{margin}(b) leaves unexplained signal at least
$C_\beta^2\varphi_0^2(1-1/n)\,\sigs^2\log n$; hence, if moreover $C_\beta^2\varphi_0^2>c_1(\sz,A)$, Assumption~\as{margin}
holds for all large $n$ with any fixed $c_{\mathrm{id}}\in(c_1(\sz,A),C_\beta^2\varphi_0^2)$.
\end{block}
{\footnotesize \textbf{PROVED}\quad Registry: \rid{W3-POST2:E}, \rid{W3-POST2:D}: \rid{PROVED-CHECKED, conditional on W2-APP:POST}.}

\end{frame}

\begin{frame}{Supplementary: the posterior (S5)}
\footnotesize
\begin{block}{Theorem S5 (Omitted variables)}
Assume \as{DIM} (with $p$ the number of candidates), Assumption~\as{omitted}, and that $s_{\mathcal C}\ge1$ is fixed. The
quantities $m_{\mathcal C}$, $\Sigma$, $\beta^{\mathcal C}$ and $\sigma^2_{\mathcal C}$ may vary with $n$; no condition beyond
Assumption~\as{omitted} is placed on how they vary. Then
\[
  \E\big[1-\pi(\Sz^{\mathcal C}\mid y)\big]=O\Big(\sqrt{\frac{\log n}{n}}\Big),\qquad
  \P\big(\HPM\neq\Sz^{\mathcal C}\big)=O\Big(\sqrt{\frac{\log n}{n}}\Big).
\]
\end{block}
{\footnotesize \textbf{PROVED}\quad Registry: \rid{W3-MISR-R:CorGb-plus}: \rid{PROVED-CHECKED, conditional on W2-APP:POST, C-MIS:GDb, C-MIS:Prop-P}; \rid{W3-MISR:MIS(i)}: \rid{PROVED-CHECKED, conditional on W2-APP:POST}.}

\end{frame}

\begin{frame}{Supplementary: the posterior (S6)}
\footnotesize
\begin{block}{Remark S6 (Smaller sets, in expectation)}
Let the raw rows be i.i.d.\ $\Normal_p(m,\Sigma)$ with $\Sigma_{jj}>0$ for every $j$, and let $R$ be the correlation
matrix of a row. Let $\yr=\alpha\one+\Xr\bstar+\varepsilon$ with $\varepsilon\sim\Normal_n(0,\sigma^2I_n)$ independent of
$\Xr$, $\sigma^2>0$, and $\alpha$, $\bstar$ deterministic; $m$, $\Sigma$, $\alpha$, $\bstar$, $\sigma^2$ and $p$ may depend on
$n$. Assume \as{DIM} and \as{SP}; $\lmin(R_{TT})\ge c$ for a fixed $c>0$ and every $T\subseteq\set p$ with
$1\le\abs{T}\le3(A+1)(\sz+1)\log n$; and $\bstar{}^\top\Sigma\bstar\le M_0\sigma^2$ for a fixed $M_0$, with
$n\min_{j\in\Sz}\Sigma_{jj}\beta^{*2}_j/(\sigma^2\log n)\to\infty$. Set $\pis{\Sz}:=0$ if $\sz>K$. Then, under the joint law
of the design and the noise, $\E[1-\pis{\Sz}]$, $\P(\HPM\neq\Sz)$, $\P(\MPM\neq\Sz)$ and
$\E\max_{j\in\set p}\abs{\PIP_j-\ind{j\in\Sz}}$ are all $O(\rn)$.
\end{block}
{\footnotesize \textbf{PROVED}\quad Registry: \rid{W3-D2:Thm2X}: \rid{PROVED-CHECKED, conditional on W2-APP:POST, C-MIS:GDb, C-MIS:Prop-P}.}

\end{frame}

\begin{frame}{Supplementary: the distribution (S7)}
\footnotesize
\begin{block}{Proposition S7 (Level-$M$ form, and the case without the collinearity condition)}
\textup{(i)} There is a constant $C_{\mathrm s}$, depending only on $(s,A,L,c_\psi,\bar b)$, such that the following holds. The
constants of Theorem~2 serve every level $M$ with $M_0\le M\le\sqrt n/\log n$, and $M$ may depend on $n$:
whenever $n\ge n_0$ and Assumption~\as{dist} holds with the window $q_n$ of this $M$,~\textup{(13)} holds, and also
\begin{equation}\tag{20}
  \P^*\Big[\,\Pi\big(\abs{S}>s+C_{\mathrm s}M\ \big|\ \Dn\big)>n^{-c_{\mathrm H} M}\Big]\ \le\ n^{-c_{\mathrm H} M}.
\end{equation}
The two exceptional sets can be taken to be the same set, which also contains every data set on which a predictor is
removed; and the bounds are uniform over all laws of the data that satisfy
Assumption~\as{dist} with the given constants.
\textup{(ii)} If Assumption~\as{dist}(d) is dropped and Assumption~\as{dist}(c) is
required only for single predictors, then~\textup{(13)} holds with radius $\sqrt{M\log n\cdot\log(\praw\vee n)/n}$,
and~\textup{(20)} with $s+C_{\mathrm s}M\log n$ in place of $s+C_{\mathrm s}M$.
\end{block}
{\footnotesize \textbf{PROVED}\quad Registry: \rid{W3-DIST:Thm-A}: \rid{PROVED-CHECKED}.}

\end{frame}

\begin{frame}{Supplementary: the distribution (S8, S9)}
\footnotesize
\begin{block}{Theorem S8 (Parametric rate when the true predictors are detectable)}
Assume Assumption~\as{dist} with the window $q_n$ of level $M=1$, with exactly $s\ge1$ nonzero coordinates of $\bstar$;
that $\lmin(\Corr(x_T))\ge\psi$ for a fixed $\psi>0$ and every set $T$ of at most $(2A+3)s+2$ predictors; and that
$\min\{\abs{\beta^*_j}\sd(x_j)/\sigs:\beta^*_j\neq0\}\ge b_n$, where $b_n>0$ and $b_n^2n/\log n\to\infty$. Then
$\Pi(d_H(f,f^*)>M_n/\sqrt n\mid\Dn)\to0$ in $\P^*$-probability, for every sequence $M_n\to\infty$.
\end{block}
{\footnotesize \textbf{PROVED}\quad Registry: \rid{W3-DIST:Thm-B}: \rid{PROVED-CHECKED}.}

\begin{block}{Remark S9 (The output of the sampler, conditionally on the mass of its list)}
Assume the conditions of Proposition~S7(i), and that the returned log scores are computed in exact
arithmetic. With $\P^*$-probability at least $1-n^{-c_{\mathrm H}M}$, every completed run, whatever its random numbers, with
$\pi(\mathcal T\mid\Dn)\ge n^{-c_{\mathrm H}M/2}$ satisfies $\hat\Pi(d_H(f,f^*)>\sqrt{M\log n/n})\le n^{-c_{\mathrm H}M/2}$. If a
completed run scored $\Sz$ and $\pi(\Sz\mid\Dn)>\tfrac12$, then the list contains $\Sz$, $\Shat=\Sz$, and $\pi(\mathcal
T\mid\Dn)>\tfrac12$. Runs that enumerate $\Om$ are not covered.
\end{block}
{\footnotesize \textbf{PROVED}\quad Registry: \rid{W3-DIST:Cor-C}: \rid{PROVED-CHECKED, conditional on EX-A (exact arithmetic)}.}

\end{frame}

\begin{frame}{Supplementary: the start law (S10)}
\footnotesize
\begin{block}{Theorem S10 (The start law concentrates)}
\textup{(i)} Assume \as{GR}, \as{DIM} with $p>\sz$, \as{FSt}, \as{V}, \as{NL}, and $\lmin(\Sigma_{n,S})\ge\varphi_{\mathrm b}^2$
for every $S$ with $\abs{S}\le16+\sz$, for a fixed $\varphi_{\mathrm b}>0$. Then $\P(\pi_0(\Sz)\ge1-n^{-c})\to1$ for every
$c<\tfrac12$. If moreover $\Cov(W_l,W_j)=0$ for every null $l$ and true $j$, then $1-\pi_0(\Sz)=\OP(\rn)$. The exponent
$\tfrac12$ cannot be improved: $1-\pi_0(\Sz)\ge\tfrac13(n+1)^{-1/2}$ for every valid data set with $n\ge4$ and $\sz<K$.

\textup{(ii)} Keep the other assumptions of~(i), and let \as{NL} hold only for the null predictors $l\notin Q$, for a
deterministic $Q\supseteq\Sz$ with $\abs{Q}\le K$ (the neighbourhoods $\Hl_l$ and the true predictors $j\in\Sz$ are
unchanged). Then $\P\big(\pi_0(\Lat(Q))\ge1-(\abs{Q}+1)n^{-1/2}e^{C\sqrt{\log n}}\big)\to1$; if $\abs{Q}\le
n^{1/2-\epsilon}$ for a fixed $\epsilon>0$, then $\P(\pi_0(\Lat(Q))\ge1-n^{-c})\to1$ for every $c<\epsilon$.
\end{block}
{\footnotesize \textbf{PROVED}\quad Registry: \rid{W3-START2-R:START-R}, \rid{W3-START2-R:BR}, \rid{W3-START:FLOOR}: \rid{PROVED-CHECKED}.}

\end{frame}

\begin{frame}{Supplementary: the reported model (S11)}
\footnotesize
\begin{block}{Theorem S11 (The reported model equals the true model: fixed design)}
Let the design be fixed, with no raw predictor constant in the sample (then $p=\praw$). Assume \as{GL}, \as{DIM}, \as{SP},
\as{BC}, \as{EF}, \as{NC}, \as{BM}(b), Assumption~\as{reported-fixed} and Convention~\as{C1}. Then the conclusions of
Theorem~3 hold: there are data events $\mathcal G_n$ with $\P(\mathcal G_n)\to1$ on which every
admissible run that produces output reports $\Shat=\Sz$, for every seed and calibration outcome; and the delivery
clauses hold under \as{C2} on $\Dom$ and under \as{C2p} for every admissible $n\ge3$.
\end{block}
{\footnotesize \textbf{PROVED}\quad Registry: \rid{W3-RM2-R:RM2-F-plus}, \rid{W3-RM2:RM2-F}: \rid{PROVED-CHECKED}.}

\end{frame}

\begin{frame}{Supplementary: the reported model (S12)}
\footnotesize
\begin{block}{Corollary S12 (Block and autoregressive designs)}
Assume \as{GR} with $\Sig$ positive definite, \as{DIM}, \as{FSt} and Convention~\as{C1}; write
$\vartheta:=(\theta^*_k)_{k\in\Sz}$ with $\Sz$ in increasing order, and let $B$ be the set of Lemma~S13 for the
design sequence. Assume that, for every $n\ge n_0$, $(\Sig,\tstar)$ belongs to one of the two families:
\begin{enumerate}[(A)]
  \item \emph{Block.} For a fixed $g\ge1$ and a fixed $c>0$, and a deterministic partition of $\set p$ into groups of size at
    most $g$, $\Sig$ is block-diagonal on the partition and $\lmin(\Sig)\ge c$; and one of: (B1) every group contains at
    most one true predictor; (Bs) for every group $G$ with $\abs{G\cap\Sz}\ge2$ there are signs $d_l\in\{-1,+1\}$, $l\in G$,
    such that $\mathrm{diag}(d)(\Sigma_{n,GG})^{-1}\mathrm{diag}(d)$ has nonpositive off-diagonal entries and
    $d_k\,\mathrm{sgn}\,\theta^*_k$ is the same for all $k\in G\cap\Sz$; (Bg) $\inf_{n\ge n_0}\mathrm{dist}(\vartheta_n,B)>0$.
  \item \emph{Autoregressive.} $\Sigma_{n,lk}=\rho^{\abs{l-k}}$ for a fixed $\rho\in(-1,1)$; and one of: (ARs)
    $\mathrm{sgn}\,\theta^*_k$ is the same for all $k\in\Sz$ if $\rho>0$, $\mathrm{sgn}\,\theta^*_k=\varsigma(-1)^k$ for a
    common $\varsigma$ if $\rho<0$, and no condition if $\rho=0$; (ARsep) $\rho\neq0$ and, for every $j\in\Sz$,
    $\sum_{k\in\Sz\setminus\{j\}}\abs{\theta^*_k}\abs\rho^{\abs{j-k}}\le\tfrac12\abs{\theta^*_j}(1-\rho^2)/(1+\rho^2)$; (ARg)
    $\inf_{n\ge n_0}\mathrm{dist}(\vartheta_n,B)>0$.
\end{enumerate}
Then the conclusions of Theorem~3 hold: there are data events $\mathcal G_n$ with $\P(\mathcal G_n)\to1$
on which every admissible run that produces output reports $\Shat=\Sz$, whatever its seed and calibration outcome; and
the delivery clauses of Theorem~3 hold under Conventions~\as{C2} and~\as{C2p}, as stated there.
\end{block}
{\footnotesize \textbf{PROVED}\quad Registry: \rid{W3-FAMC:Cor-FAMC}: \rid{PROVED-CHECKED}.}

\end{frame}

\begin{frame}{Supplementary: sufficient conditions for visibility (S13)}
\footnotesize
\begin{block}{Lemma S13 (Sufficient conditions for visibility)}
For one $n$, let $\Sig$ be a correlation matrix with $c_{16}:=\min_{\abs{V}\le16}\lmin(\Sigma_{n,V})>0$; write
$\Sz=\{k_1<\dots<k_{\sz}\}$ and $\vartheta:=(\theta^*_{k_1},\dots,\theta^*_{k_{\sz}})$. For $U\subseteq\set p$ and $j,k\notin U$
let $\Gamma^U_{jk}:=\Cov(W_j,W_k\mid W_U)$. For $j\in\Sz$ and $U\not\ni j$ with $\abs{U}\le15$, the \emph{visibility vector}
$g(U,j)\in\R^{\sz}$ has entries $\Gamma^U_{jk_i}$ for $k_i\notin U$ and $0$ for $k_i\in U$. Let
$\mathcal N_n:=\{g(U,j)/\norm{g(U,j)}\}$ over these $(U,j)$; for a design sequence, $\mathcal N$ is the closure of
$\bigcup_{n\ge n_0}\mathcal N_n$ and $B:=\bigcup_{\nu\in\mathcal N}\nu^\perp$.
\begin{enumerate}[(1)]
  \setcounter{enumi}{-1}
  \item \emph{Reduction.} $\Cov(W_j,W^\top\tstar\mid W_U)=\langle g(U,j),\vartheta\rangle$, and the entry of $g(U,j)$ at $j$ is
    $\Var(W_j\mid W_U)\ge\lmin(\Sigma_{n,U\cup\{j\}})\ge c_{16}$.
  \item \emph{Sign route.} If $\Sig$ is positive definite and there is $D=\mathrm{diag}(d_1,\dots,d_p)$, $d_l\in\{-1,+1\}$,
    such that $D\Sig^{-1}D$ has nonpositive off-diagonal entries and $d_k\,\mathrm{sgn}(\theta^*_k)=d_j\,\mathrm{sgn}(\theta^*_j)$
    for all $j,k\in\Sz$ in the same connected component of the graph of $\Sig^{-1}$, then
    $\abs{\Cov(W_j,W^\top\tstar\mid W_U)}\ge\abs{\theta^*_j}\Var(W_j\mid W_U)$ for every $j\in\Sz$ and every $U\not\ni j$, and
    this is at least $\theta_{\mathrm{lo}}c_{16}$ when $\abs{U}\le15$.
\end{enumerate}
\end{block}
\end{frame}

\begin{frame}{Supplementary: sufficient conditions for visibility (S13) (cont.)}
\footnotesize
\begin{block}{Lemma S13 (Sufficient conditions for visibility) (cont.)}
\begin{enumerate}[(1)]\setcounter{enumi}{1}
  \item \emph{Generic-position route.} For every $n\ge n_0$, $j\in\Sz$ and $U\not\ni j$ with $\abs{U}\le15$,
    $\abs{\Cov(W_j,W^\top\tstar\mid W_U)}\ge c_{16}\,\mathrm{dist}(\vartheta,B)$. The set $B$ is closed, and
    $\mathrm{dist}(\vartheta,B)=\min_{\nu\in\mathcal N}\abs{\langle\nu,\vartheta\rangle}$ (attained). Hence, if
    $\underline c:=\inf_{n\ge n_0}c_{16}(\Sig)>0$ and $\delta:=\inf_{n\ge n_0}\mathrm{dist}(\vartheta_n,B)>0$, \as{V} holds for
    every $n\ge n_0$ with $a=\underline c\,\delta$.
  \item If $\mathcal N$ is countable, $B$ is a countable union of hyperplanes through $0$: closed, Lebesgue-null and
    nowhere dense. If $\mathcal N$ is finite, the minimum in~(2) is over finitely many $\nu$.
\end{enumerate}
\end{block}
\end{frame}

\begin{frame}{Supplementary: sufficient conditions for visibility (S13) (cont.)}
\footnotesize
\begin{block}{Lemma S13 (Sufficient conditions for visibility) (cont.)}
\begin{enumerate}[(1)]\setcounter{enumi}{3}
  \item $\mathcal N$ is countable in each of the following cases, provided $\inf_{n\ge n_0}\min_{\abs{V}\le15+\sz}\lmin(\Sigma_{n,V})>0$:
    (i) the labelled matrices $\Sigma_{n,U\cup\Sz}$ take finitely many values over all $n$, $j$, $U$ (for example
    block-diagonal designs with blocks from a fixed finite list; equicorrelated blocks with a fixed correlation in $[0,1)$
    and any sizes; equicorrelation with a fixed correlation; factor models $\Sig=\Psi+\Lambda\Lambda^\top$ whose rows take
    finitely many values); (ii) stationary designs $\Sigma_{n,ll'}=r(\abs{l-l'})$ for one fixed function $r$ with
    $r(d)\to0$ as $d\to\infty$ (for example AR($\rho$) with a fixed $\abs\rho<1$, or banded designs), and block-diagonal
    designs whose blocks are of this form with finitely many such functions.
  \item \emph{Instances of the sign route}, each with its $D$ and floor: independent columns ($D=I$, $c_{16}=1$); AR($\rho$),
    $\abs\rho<1$ ($d_l=\varsigma^l$, $\varsigma:=\mathrm{sgn}\,\rho$, $\varsigma:=1$ if $\rho=0$; $c_{16}\ge(1-\abs\rho)/(1+\abs\rho)$);
    equicorrelation $\rho\in[0,1)$ ($D=I$, $c_{16}\ge1-\rho$); one factor $\Sig=\Psi+\lambda\lambda^\top$ with $\Psi$ diagonal
    positive ($d_l=\mathrm{sgn}\,\lambda_l$, $\mathrm{sgn}\,0:=1$; $c_{16}\ge\min_l\psi_l$); and block-diagonal combinations of
    these (alignment required only within a block).
\end{enumerate}
\end{block}
{\footnotesize \textbf{PROVED}\quad Registry: \rid{W3-AUDIT-V:VIS}: \rid{PROVED-CHECKED}.}

\end{frame}

\begin{frame}{Supplementary: estimates from the returned list (S14)}
\footnotesize
\begin{block}{Corollary S14 (List-renormalised estimates on the designs of Theorem~3)}
Under the assumptions of Theorem~3, on its data events $\mathcal G_n$, every admissible run that produces
output, for every seed and calibration outcome, reports $\Shat=\Sz$ and has $\Sz$ in its returned list $\mathcal T$ (on
exact enumeration, in the enumeration's list where it is returned); and then, with $\varepsilon_n:=1-\pis{\Sz}$,
\begin{gather}
  \max_j\big|\mathrm{pip}_{\mathcal T}(j)-\ind{j\in\Sz}\big|\le\varepsilon_n,\qquad
  \max_j\big|\mathrm{pip}_{\mathcal T}(j)-\PIP_j\big|\le\varepsilon_n,\tag{21}\\
  0\ \le\ \log Z-\log Z_{\mathcal T}\ \le\ -\log(1-\varepsilon_n)\quad\text{(where $\log Z_{\mathcal T}$ is reported)},\tag{22}\\
  \mathrm{mpm}_{\mathcal T}=\Shat=\Sz\quad\text{when }\varepsilon_n<\tfrac12 .\tag{23}
\end{gather}
If $p\ge2$ for all large $n$, then $\varepsilon_n=\OP(\rn)$, and for every deterministic $t_n\in(0,\tfrac12)$ with
$t_n/\rn\to\infty$ there are data events of probability tending to one on which every admissible run that produces output
has both errors in~\textup{(21)} at most $t_n$, $\mathrm{mpm}_{\mathcal T}=\Shat=\Sz$, and, where $\log Z_{\mathcal T}$ is
reported, $\abs{\log Z_{\mathcal T}-\log Z}\le2t_n$. Under Convention~\as{C2p} every run produces output, so these hold
without ``that produces output''.
\end{block}
{\footnotesize \textbf{PROVED}\quad Registry: \rid{W3-D5P:LIST-K}, \rid{W3-D5P:LIST-R}: \rid{PROVED-CHECKED}; code conformance \rid{W3-D5P:LIST-CODE}:
\rid{CODE-CHECKED (code reading at core 4676af6)}.}

\end{frame}

\begin{frame}{Supplementary: the reported estimates (S15)}
\footnotesize
In this subsection, on $\Acal$, $\Zhat$, $\PIPhat$ and $\MPMhat$ denote the pooled values the four islands would produce if
the run did not stop; on completed runs they are the reported values. Enumeration is not covered.
$\rho_n:=n^{-1/2}e^{C_\rho\sqrt{\log n}}$, with $C_\rho$ depending only on the fixed quantities of the assumptions.
\begin{block}{Remark S15 (The reported estimates are consistent)}
Assume \as{GR}, \as{DIM} with $p>\sz$, \as{FSt}, \as{EFS}, \as{V}, \as{NL}, and Conventions~\as{C1} and~\as{C3}. Then,
under the joint law of the design, the noise and the method's random numbers,
\begin{gather}
  \ind{\Acal}\Big(\big|\Zhat/Z-1\big|+\big|\log\Zhat-\log Z\big|+\max_{j\in\set{p}}\big|\PIPhat_j-\ind{j\in\Sz}\big|\Big)
    =\OP(\rho_n),\tag{24}\\
  \P\big(\Acal\cap\{\MPMhat\neq\Sz\}\big)\to0,\qquad \P\big(\text{the run completes and }\Shat\neq\Sz\big)\to0,
    \tag{25}
\end{gather}
and~\textup{(24)} also holds for the reported values multiplied by the indicator of a completed run.
\end{block}
{\footnotesize \textbf{PROVED}\quad Registry: \rid{W3-SAMP:C}, \rid{W3-SAMP:LT\#}, \rid{W3-SAMP:D}, \rid{W3-READER:R4}, \rid{W3-START2-R:IF-START}: \rid{PROVED-CHECKED}.}

\end{frame}

\begin{frame}{Supplementary: the reported estimates (S16, S17)}
\footnotesize
\begin{block}{Proposition S16 (Pooling is no worse than the worst island)}
Under Convention~\as{C3}, on $\Acal$,
$\max_j\big|\PIPhat_j-\ind{j\in\Sz}\big|\le\max_{r\le4}\max_j\big|\PIPhat_{r,j}-\ind{j\in\Sz}\big|$.
\end{block}
{\footnotesize \textbf{PROVED}\quad Registry: \rid{W3-SAMP:LT\#} (clause LT-d): \rid{PROVED-CHECKED}.}

\begin{block}{Remark S17 (An alternative route for the estimates)}
Assume Convention~\as{C3}, that $1-\pis{\Sz}\to0$ in probability, and the global-overlap condition: with
$M_\pi:=\max_{S\in\Om}\pis{S}/\pi_0(S)$, $\P(\Acal\cap\{M_\pi>\epsilon\log n\})\to0$ for every $\epsilon>0$. Then, for
every $\epsilon>0$,
\begin{gather*}
  \P\big(\Acal\cap\{\max_j\abs{\PIPhat_j-\ind{j\in\Sz}}>\epsilon\}\big)\to0,\qquad
  \P\big(\Acal\cap\{\MPMhat\neq\Sz\}\big)\to0,\\
  \P\big(\text{the run completes and }\Shat\neq\Sz\big)\to0 .
\end{gather*}
\end{block}
{\footnotesize \textbf{PROVED}\quad Registry: \rid{W3-SAMP:TR} (clause TR-d, production form): \rid{PROVED-CHECKED}. The global-overlap condition is a named
open premise (\rid{P-OVL}$\infty$); it is not derived from conditions on the design (\rid{W3-AUDIT-A:OVL} discharges only
the weaker premise \rid{P-OVL}).}

\end{frame}

\end{document}
